\section{Heat Equation}%
\label{sec:heat-equation}

The 2D heat problem on $(0,\pi)^2$ with homogeneous Dirichlet boundary conditions has the forward operator
    \[
        Kx(t_1,t_2) = \int_0^\pi \int_0^\pi k(t_1,t_2,s_1,s_2)x(s_1,s_2)\odif{s_1}\odif{s_2}
    \]
with the kernel
    \[
        k(t_1,t_2,s_1,s_2) = \sum_{n,m} e^{-(n^2+m^2)T}\sin(ns_1)\sin(ms_2)\sin(nt_1)\sin(mt_2).
    \]
It takes as input the image $x(s_1,s_2)$ as initial condition and outputs the heat distribution at time $T$, $y(t_1,t_2)
    = Kx(t_1,t_2)$.

\begin{enumerate}[label = (\alph*)]
    \item Describe in an essay text (max 10 lines) why the 2D backwards heat equation is ill-posed.
\end{enumerate}

\nopagebreak
\noindent\hrulefill

For a problem to satisfy Hadamard's well-posedness, it must meet three conditions.
These are the \emph{existence} and \emph{uniqueness} of the solution and the \emph{stability} of the problem.
The first two are trivially true for the exact heat equation.
Problem in the well-posedness of the 2D backwards heat equation lies in its stability.
Let us consider a positive error function that is corrupting our measurement of the end temperature.
The forward heat operator acts inherently as a sort of smoothing operator on the function that is at its input due to it
    being a parabolic partial differential equation problem.
This fact can be also realised by considering the growth in entropy that is instructed by the second law of
    thermodynamics.
Thus if we should invert the problem, we have to go from a smoothed function to a function that is more rigged and this
    leads to the problem with the stability of the backwards heat equation.
From another point of view, the operator norm is smaller than \(1\) hence there is intuitively some information loss
    about the original function caused by the forward operator.
Due to these aspects of the problem an arbitrarily small perturbation in the measurement temperatures may lead to
    arbitrarily large errors in the reconstructed initial temperatures.

\nopagebreak
\noindent\hrulefill

\begin{enumerate}[resume, label = (\alph*)]
    \item Implement numerically using FFT the operator $K$ as a Matlab/Python/Julia... function that takes as input
        an $N$ by $N$ matrix $x$ of sampled values of the function $x$ and returns an $N$ by $N$ matrix $y$
        approximation $y(t_1,t_2)$. Show that the function acts correct on $x(s_1,s_2) = \sin(ks_1)\sin(ls_2)$.
\end{enumerate}

\nopagebreak
\noindent\hrulefill

\begin{figure}
    \centering
    \begin{minted}[breaklines,escapeinside=||,mathescape=true, numbersep=3pt, gobble=2, frame=single,
        fontsize=\fontsize{10pt}{13pt}\selectfont, framesep=2mm]{julia1}
"""
# Arguments
- `f₀` -- initial condition
- `T` -- terminal time
"""
function heat_forward(f₀::Matrix{NT}; T::NT=0.001)::Matrix{NT} where {NT <: Number}
    @assert size(f₀, 1) == size(f₀, 2) "Expected a N×N square matrix, got $(size(f₀))"
    Nx,Ny = size(f₀)
    Kx,Ky = fftshift.((-Nx:(Nx - 1), -Ny:(Ny - 1)))
    Kx,Ky = ndgrid(Kx,Ky)
    f₀_crop = parent(padarray(f₀[2:(end-1), 2:(end-1)], Fill(zero(eltype(f₀)), (1,1), (1,1))))
    f₀_ext = [f₀_crop reverse(f₀_crop, dims=1);
            -1 .* reverse(f₀_crop, dims=2) -1 .* reverse(f₀_crop)]
    x̂ = fft(f₀_ext)
    x̂ .*= @. exp(-(Kx^2 + Ky^2) * T)
    y = ifft(x̂)
    y = real(y[1:Nx, 1:Ny])
    return y
end
    \end{minted} 
    \caption{A \emph{Julia} implementation of the 2D forward heatoperator using FFT. The used packages are \texttt{LazyGrids.jl},
        \texttt{ImageFiltering.jl} (for the \texttt{padarray} function) and \texttt{FFTW.jl}.}
    \label{fig:julia-implementation}
\end{figure}

The \emph{Julia} implementation of the operator is in \cref{fig:julia-implementation}.

The exact analytical solution of the  \(Kx\) when \(x(s_1, s_2) = \sin(ks_1) \sin(ks_2)\) is \[
    y(t_1, t_2) = e^{-2T^2} \sin(ks_1) \sin(ks_2),
    \] thus it should effectively only lower the modes of the harmonics since the sine functions are eigenfunctions for the
    operator.

A test of the implementation is made using the harmonic eigenfunctions in \cref{fig:test-heat-fft}.

\begin{figure}
    \centering
    \begin{subfigure}[t]{0.48\textwidth}
        \includegraphics[width = \textwidth]{fig_initial_condition.pdf}
        \caption{The initial function \(x(4,3)\).}
        \label{fig:initial-heat}
    \end{subfigure}
    \begin{subfigure}[t]{0.48\textwidth}
        \includegraphics[width = \textwidth]{fig_forward_solution.pdf}
        \caption{Solution of the forward heat equation \(K(x(4,3))\) with the terminal time \(T = 0.02\).}
        \label{fig:forward-heat}
    \end{subfigure}
    \caption{A test of the implementation of the FFT forward heat equation operator using harmonic eigenfunctions.}
    \label{fig:test-heat-fft}
\end{figure}

\nopagebreak
\noindent\hrulefill

\begin{enumerate}[resume, label = (\alph*)]
    \item Download the image \texttt{KarinHarbo.jpg} and create from pixels in the image a 512 by 512 matrix $x$. Fix $T
        = 0.001$ and compute $y = Kx$. Add random Gaussian noise to $y$ of magnitude $\delta = 0.01$ to find
        $y^\delta$. Display $y^\delta$.
\end{enumerate}

\nopagebreak
\noindent\hrulefill

The field created from the image (\cref{fig:karin-harbo-field}), forward passed image (\cref{fig:karin-harbo-forward})
    and the forward passed image with the added noise (\cref{fig:karin-harbo-noisy}).
The image was converted to grayscale, resized to \(512\times 512\) pixels and the edges where mirrored to generate
    a border of \(30\) pixels around all of the image borders and tapered by the cosine apodization function to ensure
    that we comply with the condition of the array being periodic\footnote{Edges of the array are gradually (in the span
    or \(30\) pixels) attenuated to \(0\) at both sides, hence the array concatenated with itself at all the sides is
    periodic in compliance with the assumptions.} as expected by the FFT algorithm\footnote{This is done to alleviate
    artifacts that can be caused by the FFT without its assumptions satisfied.}.
You can see that with this noise level, the added white noise is not even visible in \cref{fig:karin-harbo}.
\begin{figure}
    \centering
    \begin{subfigure}[t]{0.48\textwidth}
        \includegraphics[width = \textwidth]{fig_karin_harbo.pdf}
        \caption{Karin Harbo image field \(x\)}
        \label{fig:karin-harbo-field}
    \end{subfigure}
    \hfill
    \begin{subfigure}[t]{0.48\textwidth}
        \includegraphics[width = \textwidth]{fig_karin_harbo_forward.pdf}
        \caption{Forward passed Karin Harbo image \(Kx = y\)}
        \label{fig:karin-harbo-forward}
    \end{subfigure}
    \begin{subfigure}[t]{0.48\textwidth}
        \includegraphics[width = \textwidth]{fig_karin_harbo_forward_noisy.pdf}
        \caption{Forward passed Karin Harbo image with added noise of magnitude \(10\%\), \(y^{\delta}\)}
        \label{fig:karin-harbo-noisy}
    \end{subfigure}
    \hfill
    \begin{subfigure}[t]{0.48\textwidth}
        \includegraphics[width = \textwidth]{fig_landweber_reconstruction.pdf}
        \caption{Reconstruction of the Karin Harbo image from the heated and noisy \(y^{\delta}\) using the Landweber
        iterations with the relaxation parameter \(\omega = 1\) and \(500\) iterations}
        \label{fig:karin-harbo-landweber}
    \end{subfigure}
    \caption{Karin Harbo original image, forward passed and with added noise.}
    \label{fig:karin-harbo}
\end{figure}
    
\nopagebreak
\noindent\hrulefill

\begin{enumerate}[resume, label = (\alph*)]
    \item Show that $K$ is self-adjoint. Solve the inverse problem using the Landweber iteration (take 500 iterations).
\end{enumerate}

\nopagebreak
\noindent\hrulefill

\paragraph{Self-adjoint:}%
\label{par:self-adjoint}

An operator is self-adjoint if \[
    \langle Ku, v \rangle = \langle u, Kv \rangle, \quad \forall x,y \in L^2(\Omega).
\] 
In our case, we have \(\Omega = \left(0, \pi  \right)^2\).
We can write the first inner product as \[
    \langle Ku, v \rangle = \iint_{\left(0, \pi  \right)^2} \overline{v(t_1, t_2)} \iint_{\left(0, \pi  \right)^2}
    k(t_1, t_2, s_1, s_2) u(s_1, s_2) \odif{{ (s_1,s_2) }}\odif{{ (t_1,t_2) }},
\] and similarly \[
    \langle u, Kv \rangle = \iint_{\left(0, \pi  \right)^2} \overline{u(s_1, s_2)} \iint_{\left(0, \pi  \right)^2}
    k(s_1, s_2,t_1, t_2) v(t_1, t_2) \odif{{ (t_1,t_2) }}\odif{{ (s_1,s_2) }}.
\] 

First of all, we are dealing only with real functions hence we can drop the conjucations.
Now we see that for these to be equal, we only need to check that the kernel \(k\) is symmetric, that is \[
    k(t_1, t_2, s_1, s_2) = k(s_1, s_2, t_1, t_2).
\] 

For the kernel \[
    k(t_1,t_2,s_1,s_2) = \sum_{n,m} e^{-(n^2+m^2)T}\sin(ns_1)\sin(ms_2)\sin(nt_1)\sin(mt_2),
\] we immediately see that this is true so the operator is self-adjoint.

\paragraph{Landweber iteration:}%
\label{par:landweber-iteration}

Landweber iteration is a method for sequentially approximating the inverse problem solution. We start with \(x^{(0)}
    \equiv 0 \) and iterate using the prescription
    \[
        x^{\left(k+1  \right)} = x^{\left(k  \right)} + \omega K^{*}\left(y^{\delta} - K(x^{\left(k  \right)})  \right),
    \] where \(\omega\) is the relaxation parameter of the method.
For \(|\omega| < 1/\|K^{*}K\|\) we have a guarantee of convergence.
Since for the heat operator it holds that \(\|K\| \le 1\), we can drop this parameter and the convergence is
    guaranteed\footnote{However the convergence rate is not necessarily fast.}.

From the paragraph \ref{par:self-adjoint} above we know that the operator is self-adjoint so the method changes to \[
        x^{\left(k+1  \right)} = x^{\left(k  \right)} + K\left(y^{\delta} - K(x^{\left(k  \right)})  \right).
\] 

You can see the Landweber reconstruction of the Karin Harbo image in \cref{fig:karin-harbo-landweber}.
The intermediate results of the iteration are visible in \cref{fig:karin-harbo-landweber-intermediate}.

\begin{figure}
  \centering
  \includegraphics[width=\textwidth]{fig_landweber_intermediates.pdf}
  \caption{Intermediate results of the Landweber reconstruction of the Karin Harbo image along with the change in
    relative error from the previous logged iteration. The rate of convergence slowes down significantly after the first
    few iterations as can be expected.}
  \label{fig:karin-harbo-landweber-intermediate}
\end{figure}

\nopagebreak
\noindent\hrulefill

\section{Abel Problem}%
\label{sec:abel-problem}

Consider the Abel problem for $x(s) = 1/2$, $y(t) = \sqrt{t}$. Discretize $K$ by a matrix $A \in \mathbb{R}^{256\times
    256}$ and discretize similarly $x(s)$, $y(t)$ by $x,y \in \mathbb{R}^{256}$. Add random noise of magnitude 10\% to
        $y$ to get $y^\delta$. Denote by $L_1$ a finite difference matrix approximating the derivative.

\nopagebreak
\noindent\hrulefill

A plot of the objects from the Abel problem is in \cref{fig:abel-init}.

\begin{figure}
  \centering
  \includegraphics[width=\textwidth]{fig_abel.pdf}
  \caption{The true solution \(y\), forward solution \(Kx\) and the noisy \(y^\delta\).}
  \label{fig:abel-init}
\end{figure}

\nopagebreak
\noindent\hrulefill

\begin{enumerate}[label = (\alph*)]
    \item Show that $x \in H^1(0,\pi)$. Thus it makes sense to apply $H^1$ regularization.
\end{enumerate}

\nopagebreak
\noindent\hrulefill

In order for \(x\) to belong to the Sobolev space \(H^{1}\left(0,\pi  \right)\), we need to check the it belongs to
    \(L^2(0,\pi)\), which is evident due to the boundedness of both the function \(x\) and the integration set. 
Additionally we need to check that the weak derivative \(x' \in L^2(0, \pi)\), which is also evident since it has the
    strong derivative equal to \(0\).

\nopagebreak
\noindent\hrulefill


\begin{enumerate}[resume, label = (\alph*)]
    \item We approximate the $H^1$-norm by $\|x\|^2_2 + \|L_1x\|^2_2$. 

    Show that the problem
        \begin{equation}
            \label{eq:abel-tikhonov}
            x_\alpha = \argmin_x \|Ax-y^\delta\|^2_2 + \alpha(\|x\|^2_2 + \|L_1x\|^2_2)
        \end{equation}
    has a unique solution for any $\alpha > 0$.
\end{enumerate}

\nopagebreak
\noindent\hrulefill

We can write the solution of \(x_\alpha\) equivalently as \[
    x_\alpha = \argmin_x \left\| \begin{pmatrix}
        A \\ \sqrt{\alpha} L 
    \end{pmatrix} x - \begin{pmatrix}
        y^\delta \\ 0
    \end{pmatrix} \right\|_2,
\]  similarly to the Tikhonov solution.
Here the matrix \(L\) is now composed from the Sobolev discretized regularization matrix \( L_\text{Tikhonov} = I \)
and the approximating matrix of the derivative \(L_1\) by concatenating them vertically as \[
    L = \begin{pmatrix}
        L_\text{Tikhonov} \\ L_1
    \end{pmatrix}.
\] 
Now this is a least squares problem that can be solved as \[
    x_\alpha = y^{\delta} \backslash R \quad \text{where} \quad R = \begin{pmatrix}
        A \\ \sqrt{\alpha}L
    \end{pmatrix}.
\] 
This problem has a unique solution if the matrix \(R\) has full column rank, which it obviously has.

\nopagebreak
\noindent\hrulefill
    
\begin{enumerate}[resume, label = (\alph*)]
    \item Find the solution to the ordinary Tikhonov problem (without $L_1$) with $\alpha$ obtained by the discrepancy
        principle.
\end{enumerate}

\nopagebreak
\noindent\hrulefill

The discrepancy principle states that there is a unique \(\alpha\) for which \(F(\alpha) \coloneqq \| K x_\alpha
- y^{\delta}\| = \delta\) and this \(\alpha\) is admissible, i.e. \(\lim_{\delta \to 0} x_\alpha = x\).

We find the \(\alpha\) from the discrepancy by evaluating \(F\) at various \(\alpha\) and solving the equation
    \(F(\alpha) = \delta\) numerically by the optimization of \(|F(\alpha) - \delta|\).
A graph showing \(F\) dependency on \(\alpha\) is in \cref{fig:discrepancy-tikhonov}.
\(\alpha\) chosen by the discrepancy principle turns out to be \(8.529966\times {10}^{-4}\).
For this \(\alpha\) the regularized solution is is in \cref{fig:abel-tikhonov-solution}.

\begin{figure}
  \centering
  \begin{subfigure}[t]{\textwidth}
    \includegraphics[width=\textwidth]{fig_discrepancy_tikhonov.pdf}
    \caption{\(F(\alpha)\) from the discrepancy for the Tikhonov regularization.}
    \label{fig:discrepancy-tikhonov}
  \end{subfigure}
  \begin{subfigure}[t]{\textwidth}
    \includegraphics[width=\textwidth]{fig_abel_tikhonov.pdf}
    \caption{The solution of the Tikhonov regularized inverse problem using the \(\alpha\) chosen by the discrepancy
        principle.}
    \label{fig:abel-tikhonov-solution}
  \end{subfigure}
  \caption{Tikhonov approach to the regularization of the Abel problem.}
  \label{fig:abel-tikhonov}
\end{figure}

\nopagebreak
\noindent\hrulefill
    
\begin{enumerate}[resume, label = (\alph*)]
    \item Find the solution to the problem in \cref{eq:abel-tikhonov}, again choosing $\alpha$ by the discrepancy
        principle. Compare this solution to the previous without $L_1$.
\end{enumerate}

We find the \(\alpha\) from the discrepancy principle analogously to the previous case.
Now the \(\alpha\) chosen by the discrepancy is \(2.682426\times 10^{-4}\).
Both the solution and the dependency of \(F\) on \(\alpha\) are shown in \cref{fig:abel-sobolev}.

\begin{figure}
  \centering
  \begin{subfigure}[t]{\textwidth}
    \includegraphics[width=\textwidth]{fig_discrepancy_sobolev.pdf}
    \caption{\(F(\alpha)\) from the discrepancy for the Sobolev regularization.}
    \label{fig:discrepancy-sobolev}
  \end{subfigure}
  \begin{subfigure}[t]{\textwidth}
    \includegraphics[width=\textwidth]{fig_abel_sobolev.pdf}
    \caption{The solution of the Sobolev regularized inverse problem using the \(\alpha\) chosen by the discrepancy
        principle.}
    \label{fig:abel-sobolev-solution}
  \end{subfigure}
  \caption{Sobolev approach to the regularization of the Abel problem.}
  \label{fig:abel-sobolev}
\end{figure}

A comparison of the solution for the Tikhonov and the Sobolev regularization is in \cref{fig:abel-comparison}.

The results point to three problems,
\begin{enumerate}[label = \roman*)]
    \item The Tikhonov regularization is not appropriate for the Abel problem when we expect that sum of the solution is
        not close to \(0\).
    \item Despite the admissibility of the \(\alpha\) chosen by the discrepancy principle, it is not guaranteed that for
        a noise power that is not close to \(0\) the parameter is chosen optimally.
    \item Since \(\|x\|\) is expected to be much greater than \(\|L_1 x\|\), the Sobolev norm is mainly lead by the
        ordinary \(L_2\) norm, hence the similarity in the two solutions visible from the comparison in
        \cref{fig:abel-comparison}.
\end{enumerate}
    
\begin{figure}
  \centering
  \begin{subfigure}[t]{\textwidth}
    \includegraphics[width=\textwidth]{fig_diff_tikhonov_sobolev.pdf}
    \caption{The difference of the two solutions using different regularizations.}
    \label{fig:diff-tikhonov-sobolev}
  \end{subfigure}
  \begin{subfigure}[t]{\textwidth}
    \includegraphics[width=\textwidth]{fig_comparison_tikhonov_sobolev.pdf}
    \caption{Comparison of the solutions using the Tikhonov regularization \(x_{\alpha}\) and the Sobolev regularization
    \(x_{\alpha}^{(L_2, L_1)}\)}
    \label{fig:comparison-tikhonov-sobolev}
  \end{subfigure}
  \caption{A comparison of both approaches to the regularization of the Abel problem using the parameters \(\alpha\)
    determined by the discrepancy principle.}
  \label{fig:abel-comparison}
\end{figure}

Additionally I tried to use only the second term from the Sobolev norm (which is in theory similar to the total
    variation) for the regularization i.e. finding \[
    x_{\alpha}^{(L_1)} = \argmin_{x} \|A x - y^{\delta}\|_2^2 + \alpha \|L_1 x\|_2^2.
\] 
Since we are looking for \(x\) that has a derivative of \(0\) throughout the whole domain, this should be an ideal
    regularization.
Using \(\alpha\) chosen by the discrepancy principle (\(\alpha = 3.689282 \times 10^{-4}\)) does not work very well, due
    to the reasons stated above.
I chose a large \(\alpha\) of  \(0.65\).
The resulting reconstruction may be seen in \cref{fig:abel-l1}.
Though it does improve upon the previous approaches, it still contains much to be desired.
For instance the oscillatory jumps around the solution suggest that the \(L_1\) matrix is not a great fit for estimating
    the derivative of a noisy function since the noise is not correlated.
We could improve this by using an approximation matrix derived from a higher order Runge-Kutta method approximation.
The reason for this theorised improvement is, that more terms will play a role in the estimate and hence the white noise
    will be partially removed by this term averaging. 

\begin{figure}
  \centering
  \includegraphics[width=\textwidth]{fig_abel_l1.pdf}
  \caption{The solution of the Abel problem using the regularization by the \(L_1\) matrix, \(x_{\alpha}^{(L_1)}\).}
  \label{fig:abel-l1}
\end{figure}

\nopagebreak
\noindent\hrulefill
